\documentclass{article}
\usepackage{amsmath}
\usepackage{amssymb}
\usepackage{listings}
\usepackage{xcolor}

\definecolor{codegreen}{rgb}{0,0.6,0}
\definecolor{codegray}{rgb}{0.5,0.5,0.5}
\definecolor{codepurple}{rgb}{0.58,0,0.82}

\lstdefinestyle{mystyle}{
    commentstyle=\color{codegreen},
    keywordstyle=\color{magenta},
    numberstyle=\tiny\color{codegray},
    stringstyle=\color{codepurple},
    basicstyle=\ttfamily\footnotesize,
    breakatwhitespace=false,
    breaklines=true,
    captionpos=b,
    keepspaces=true,
    numbers=left,
    numbersep=5pt,
    showspaces=false,
    showstringspaces=false,
    showtabs=false,
    tabsize=2
}

\lstset{style=mystyle}

\title{TeX Macro Compiler Tool}
\author{Autonomous AI Polyglot Software Engineer}
\date{\today}

\begin{document}

\maketitle

\section{Introduction}

This document describes the \texttt{TeX Macro Compiler Tool}, a specialized tool for optimizing macro expansions in TeX documents. The compiler is designed to handle complex document generation by efficiently processing macro definitions.

\section{Installation}

To use the \texttt{TeX Macro Compiler Tool}, ensure you have a TeX distribution installed on your system. The compiler itself is a single \texttt{.tex} file that can be executed using any standard TeX engine (e.g., \texttt{pdflatex}, \texttt{xelatex}).

\begin{lstlisting}[language=TeX]
% Execute the compiler with the following command:
% pdflatex main.tex
\end{lstlisting}

\section{Usage}

The \texttt{TeX Macro Compiler Tool} can be used to optimize macro expansions in your TeX documents. Simply provide the input file containing your macro definitions, and the compiler will generate an optimized output file.

\begin{lstlisting}[language=TeX]
\input{macros.tex} % Input macro definitions
\processmacros     % Process and optimize macros
\outputresult      % Output the optimized result
\end{lstlisting}

\section{Example}

Below is an example of a TeX macro definition and its optimized output:

\subsection{Input Macro Definition}

\begin{lstlisting}[language=TeX]
\def\greeting#1{Hello, #1!}
\def\farewell#1{Goodbye, #1!}
\def\message{\greeting{Alice} \farewell{Bob}}
\end{lstlisting}

\subsection{Optimized Output}

\begin{lstlisting}[language=TeX]
\def\greeting#1{Hello, #1!}
\def\farewell#1{Goodbye, #1!}
\def\message{Hello, Alice! Goodbye, Bob!}
\end{lstlisting}

\section{Conclusion}

The \texttt{TeX Macro Compiler Tool} provides a powerful way to optimize macro expansions in TeX documents. By reducing the number of expansion steps, it can significantly improve the performance of complex document generation workflows.

\section{Author}

This tool was created by the Autonomous AI Polyglot Software Engineer.

\end{document}
